\documentclass[11pt,a4paper]{article}

% --- Packages Standard ---
\usepackage[utf8]{inputenc}
\usepackage[T1]{fontenc}
\usepackage[french]{babel}
\usepackage{geometry}
\geometry{top=2.5cm, bottom=2.5cm, left=2.5cm, right=2.5cm}
\usepackage{graphicx}
\usepackage{xcolor}
\usepackage{hyperref}
\usepackage{titlesec}
\usepackage{fancyhdr}
\usepackage{booktabs}
\usepackage{amsmath}
\usepackage{amssymb}
\usepackage{enumitem}
\usepackage{listings}
\usepackage{caption}
\usepackage{float}

% --- Couleurs Personnalisées ---
\definecolor{navyblue}{HTML}{1E3A8A}
\definecolor{teal}{HTML}{0D9488}
\definecolor{darkgray}{HTML}{374151}
\definecolor{lightgray}{HTML}{F3F4F6}
\definecolor{codebackground}{HTML}{F9FAFB}

% --- Configuration de Hyperref ---
\hypersetup{
    colorlinks=true,
    linkcolor=navyblue,
    filecolor=teal,      
    urlcolor=teal,
    citecolor=navyblue,
    pdftitle={Rapport de Projet - Préparation Métier},
    pdfpagemode=FullScreen,
}

% --- Configuration des Titres ---
\titleformat{\section}
  {\color{navyblue}\normalfont\Large\bfseries}
  {\thesection}{1em}{}[{\titlerule[0.8pt]}]
\titleformat{\subsection}
  {\color{teal}\normalfont\large\bfseries}
  {\thesubsection}{1em}{}

% --- Configuration En-têtes et Pieds de Page (Style SUP'COM) ---
\pagestyle{fancy}
\fancyhf{}
\fancyhead[R]{\textcolor{darkgray}{\small Projet de Préparation Métier}}
\fancyfoot[C]{\thepage}
\renewcommand{\headrulewidth}{0.5pt}
\renewcommand{\footrulewidth}{0pt}

% --- Configuration du Code ---
\lstset{
    backgroundcolor=\color{codebackground},
    basicstyle=\ttfamily\small\color{darkgray},
    breakatwhitespace=false,         
    breaklines=true,                 
    captionpos=b,                    
    keepspaces=true,                 
    showspaces=false,                
    showstringspaces=false,
    showtabs=false,                  
    tabsize=2,
    frame=single,
    rulecolor=\color{lightgray}
}

\begin{document}

% ==========================================
% PAGE 1 : PAGE DE GARDE (Style SUP'COM P2M exact)
% ==========================================
\begin{titlepage}
    \centering
    \vspace*{-1cm}
    
    % Logo et En-tête Institutionnel
    {\large\textbf{\color{navyblue}{SUP'COM}}} \\
    \vspace{0.2cm}
    {\small\textbf{Ecole Supérieure des Communications de Tunis}} \\
    {\small Higher School of Communications of Tunis} \\
    \vspace{0.3cm}
    \rule{6cm}{1pt} \\
    
    \vspace{2.5cm}
    
    \rule{\linewidth}{1.5pt} \\[0.4cm]
    {\Huge\bfseries\color{navyblue}{Projet de Préparation Métier}} \\
    \vspace{0.2cm}
    {\large\bfseries\color{darkgray}{P2M}} \\
    \rule{\linewidth}{1.5pt} \\
    
    \vspace{2cm}
    
    {\Large\bfseries\color{darkgray}{[Saisir le titre du projet ici]}} \\[0.3cm]
    {\large\itshape\color{teal}{[Sous-titre ou description du projet]}} \\
    
    \vspace{3.5cm}
    
    \begin{minipage}[t]{0.45\textwidth}
        \flushleft
        \textbf{Réalisé par :}\\
        {\large\color{navyblue}{[Prénom Nom Étudiant 1]}} \\
        {\large\color{navyblue}{[Prénom Nom Étudiant 2]}}
    \end{minipage}
    \hfill
    \begin{minipage}[t]{0.45\textwidth}
        \flushright
        \textbf{Encadré par :}\\
        {\large\color{navyblue}{[Prénom Nom de l'Encadrant]}}
    \end{minipage}
    
    \vfill
    
    \rule{8cm}{0.4pt} \\
    \vspace{0.3cm}
    {\large\textbf{Année universitaire : 2025--2026}}
\end{titlepage}

\newpage
\pagenumbering{roman}
\setcounter{page}{1}

% ==========================================
% PAGE i : REMERCIEMENTS
% ==========================================
\section*{Remerciements}
\addcontentsline{toc}{section}{Remerciements}
\textit{[Saisir votre texte de remerciements ici. Remerciez votre encadrant, les membres du jury et toute personne ayant contribué au projet.]}

\newpage

% ==========================================
% PAGE ii : TABLE DES MATIÈRES
% ==========================================
\renewcommand{\contentsname}{Table des matières}
\tableofcontents
\addcontentsline{toc}{section}{Table des matières}

\newpage

% ==========================================
% PAGE iii : LISTE DES FIGURES
% ==========================================
\renewcommand{\listfigurename}{Liste des figures}
\listoffigures
\addcontentsline{toc}{section}{Liste des figures}

\newpage

% ==========================================
% PAGE iv : LISTE DES ABRÉVIATIONS
% ==========================================
\section*{Liste des abréviations}
\addcontentsline{toc}{section}{Liste des abréviations}

\begin{description}[leftmargin=2.5cm, style=nextline]
    \item[EMS] Element Management System
    \item[RTU] Remote Test Unit
    \item[OTDR] Optical Time-Domain Reflectometer
    \item[OTH] Optical Test Head
    \item[OSM] OpenStreetMap
    \item[GIS] Geographic Information System
    \item[[Abreviation]] [Signification de l'abréviation]
\end{description}

\newpage
\pagenumbering{arabic}
\setcounter{page}{1}

% ==========================================
% SECTION 1 : INTRODUCTION
% ==========================================
\section{Introduction}
\textit{[Saisir l'introduction générale du projet ici. Présentez brièvement le projet de Préparation Métiers (P2M), son objectif global et les grandes lignes du travail.]}

\subsection{État de l’art}

\subsubsection{Contexte}
\textit{[Saisir le contexte général ici. Expliquez le domaine des réseaux de fibres optiques, leur importance, leur utilisation dans les infrastructures modernes de télécommunications.]}

\subsubsection{Problématique}
\textit{[Saisir la problématique ici. Décrivez les limites des méthodes traditionnelles, la nature imprévisible des coupures physiques de fibres optiques, et l'impact opérationnel ou financier des pannes.]}

\subsubsection{Solution}
\textit{[Saisir la solution proposée ici. Détaillez l'application web ou le système conçu pour répondre aux besoins opérationnels de supervision et de maintenance.]}

\subsection{Conclusion}
\textit{[Saisir la conclusion de l'introduction ici.]}

\newpage

% ==========================================
% SECTION 2 : CONCEPTION
% ==========================================
\section{Conception}
\textit{[Saisir l'introduction de la section Conception ici. Présentez la démarche de conception adoptée.]}

\subsection{Architecture d’un réseau de fibre optique}
\textit{[Décrivez l'architecture générale de supervision. Expliquez le rôle des différents composants (EMS, RTU, OTH, OTDR).]}

\begin{figure}[H]
    \centering
    % Remplacer "placeholder_image" par votre fichier d'image réel
    % \includegraphics[width=0.85\textwidth]{placeholder_image.png}
    \fbox{\begin{minipage}{0.85\textwidth}
        \centering
        \vspace{1.5cm}
        \textbf{[FIGURE 2.1 : Architecture typique de supervision d'un réseau de fibre optique]} \\
        \textit{(Insérez ici votre schéma d'architecture)} \\
        \vspace{1.5cm}
    \end{minipage}}
    \caption{Architecture typique de supervision d’un réseau de fibre optique.}
    \label{fig:typical_arch}
\end{figure}

\subsection{Outils et bibliothèques}
\textit{[Présentez brièvement les outils technologiques adoptés pour la réalisation de l'application.]}

\subsubsection{Django}
\textit{[Saisir la description de l'outil et son utilité dans le projet ici.]}

\subsubsection{PostgreSQL}
\textit{[Saisir la description de l'outil et son utilité dans le projet ici.]}

\subsubsection{Node.js}
\textit{[Saisir la description de l'outil et son utilité dans le projet ici.]}

\subsubsection{Firebase}
\textit{[Saisir la description de l'outil et son utilité dans le projet ici.]}

\subsubsection{OpenStreetMap}
\textit{[Saisir la description de l'outil et son utilité dans le projet ici.]}

\subsection{Conclusion}
\textit{[Saisir la conclusion de la partie conception.]}

\newpage

% ==========================================
% SECTION 3 : IMPLÉMENTATIONS ET RÉSULTATS
% ==========================================
\section{Implémentations et résultats}
\textit{[Présentez l'introduction de la partie implémentation, montrant le passage de la conception à la réalisation concrète.]}

\subsection{Gestion des utilisateurs et des routes optiques}
\textit{[Saisir le texte ici. Présentez les interfaces d'authentification, de configuration des sondes RTU et des commutateurs OTH.]}

\begin{figure}[H]
    \centering
    \fbox{\begin{minipage}{0.85\textwidth}
        \centering
        \vspace{1.5cm}
        \textbf{[FIGURE 3.1 : Capture d'écran - Interface d'authentification]} \\
        \vspace{1.5cm}
    \end{minipage}}
    \caption{Interface d'authentification des utilisateurs.}
    \label{fig:auth}
\end{figure}

\subsection{Communication EMS/RTU}
\textit{[Saisir le texte ici. Expliquez le fonctionnement de la communication entre l'émulateur RTU (envoi des fichiers de mesure .sor) et le système centralisé EMS.]}

\subsection{Mise en action}
\textit{[Saisir le texte ici. Décrivez le scénario de détection d'une panne, l'affichage de la courbe OTDR de défaut, et le traitement cartographique cartographique GIS.]}

\begin{figure}[H]
    \centering
    \fbox{\begin{minipage}{0.85\textwidth}
        \centering
        \vspace{1.5cm}
        \textbf{[FIGURE 3.2 : Capture d'écran - Visualisation cartographique des liaisons]} \\
        \vspace{1.5cm}
    \end{minipage}}
    \caption{Visualisation cartographique des liaisons optiques et localisation d'un défaut.}
    \label{fig:map}
\end{figure}

\subsection{Conclusion}
\textit{[Saisir la conclusion de la partie réalisation.]}

\newpage

% ==========================================
% SECTION 4 : CONCLUSION ET PERSPECTIVES
% ==========================================
\section{Conclusion et Perspectives}
\textit{[Saisir la conclusion générale ici. Rappelez les objectifs atteints, le bilan personnel et technique du projet, ainsi que les perspectives d'évolutions futures.]}

\newpage

% ==========================================
% SECTION 5 : BIBLIOGRAPHIE
% ==========================================
\section*{Bibliographie}
\addcontentsline{toc}{section}{Bibliographie}

\begin{enumerate}[label={[\arabic*]}]
    \item Site Web de EXFO -- \textit{EXpert in Fiber Optics}, consulté le 26 mai 2025. \\
    Disponible sur : \url{https://www.exfo.com/en/}
    
    \item Viavi Solutions -- \textit{Outils de test et de diagnostic pour réseaux fibre optique}, consulté le 26 mai 2025. \\
    Disponible sur : \url{https://www.viavisolutions.com}
    
    \item \textit{[Saisir d'autres références de livres, sites ou articles scientifiques ici...]}
\end{enumerate}

\end{document}
